%% kapa: exam=1 round=31 year=2020 period=2 form=A subject=accounting subject_ko=회계학 seq=9 number=49
%% source: Q-Net 감정평가사 기출문제 (cst00309) artlSeq=5207219 file=extracted/1008/5207219/A형/A형/2020년도 제31회 감정평가사 1차 시험_2교시_A형.pdf
%% conversion: kordoc 4.24.51 -> tools/gen1.py (auto); review=proofread-sonnet; flags=none
\begin{question}{49}
상품매매기업인 (주)감평은 계속기록법과 실지재고조사법을 병행하고 있다. (주)감평의 20x1년 기초재고는 ￦10,000(단가 ￦100)이고, 당기매입액은 ￦30,000(단가 ￦100), 20x1년 말 현재 장부상 재고수량은 70개이다. (주)감평이 보유하고 있는 재고자산은 진부화로 인해 단위당 순실현가능가치가 ￦80으로 하락하였다.
(주)감평이 포괄손익계산서에 매출원가로 ￦36,000을 인식하였다면, (주)감평의 20x1년 말 현재 실제재고수량은? (단, 재고자산감모손실과 재고자산평가손실은 모두 매출원가에 포함한다.) \par
\begin{choices}
  \item 40개
  \item 50개
  \item 65개
  \item 70개
  \item 80개
\end{choices}
\end{question}
\answer{2}
